% Options for packages loaded elsewhere
% Options for packages loaded elsewhere
\PassOptionsToPackage{unicode}{hyperref}
\PassOptionsToPackage{hyphens}{url}
%
\documentclass[
  11pt,
  ignorenonframetext,
  aspectratio=169,
]{beamer}
\newif\ifbibliography
\usepackage{pgfpages}
\setbeamertemplate{caption}[numbered]
\setbeamertemplate{caption label separator}{: }
\setbeamercolor{caption name}{fg=normal text.fg}
\beamertemplatenavigationsymbolsempty
% remove section numbering
\setbeamertemplate{part page}{
  \centering
  \begin{beamercolorbox}[sep=16pt,center]{part title}
    \usebeamerfont{part title}\insertpart\par
  \end{beamercolorbox}
}
\setbeamertemplate{section page}{
  \centering
  \begin{beamercolorbox}[sep=12pt,center]{section title}
    \usebeamerfont{section title}\insertsection\par
  \end{beamercolorbox}
}
\setbeamertemplate{subsection page}{
  \centering
  \begin{beamercolorbox}[sep=8pt,center]{subsection title}
    \usebeamerfont{subsection title}\insertsubsection\par
  \end{beamercolorbox}
}
% Prevent slide breaks in the middle of a paragraph
\widowpenalties 1 10000
\raggedbottom
\AtBeginPart{
  \frame{\partpage}
}
\AtBeginSection{
  \ifbibliography
  \else
    \frame{\sectionpage}
  \fi
}
\AtBeginSubsection{
  \frame{\subsectionpage}
}
\usepackage{iftex}
\ifPDFTeX
  \usepackage[T1]{fontenc}
  \usepackage[utf8]{inputenc}
  \usepackage{textcomp} % provide euro and other symbols
\else % if luatex or xetex
  \usepackage{unicode-math} % this also loads fontspec
  \defaultfontfeatures{Scale=MatchLowercase}
  \defaultfontfeatures[\rmfamily]{Ligatures=TeX,Scale=1}
\fi
\usepackage{lmodern}

\ifPDFTeX\else
  % xetex/luatex font selection
\fi
% Use upquote if available, for straight quotes in verbatim environments
\IfFileExists{upquote.sty}{\usepackage{upquote}}{}
\IfFileExists{microtype.sty}{% use microtype if available
  \usepackage[]{microtype}
  \UseMicrotypeSet[protrusion]{basicmath} % disable protrusion for tt fonts
}{}
\makeatletter
\@ifundefined{KOMAClassName}{% if non-KOMA class
  \IfFileExists{parskip.sty}{%
    \usepackage{parskip}
  }{% else
    \setlength{\parindent}{0pt}
    \setlength{\parskip}{6pt plus 2pt minus 1pt}}
}{% if KOMA class
  \KOMAoptions{parskip=half}}
\makeatother

\usepackage{color}
\usepackage{fancyvrb}
\newcommand{\VerbBar}{|}
\newcommand{\VERB}{\Verb[commandchars=\\\{\}]}
\DefineVerbatimEnvironment{Highlighting}{Verbatim}{commandchars=\\\{\}}
% Add ',fontsize=\small' for more characters per line
\usepackage{framed}
\definecolor{shadecolor}{RGB}{241,243,245}
\newenvironment{Shaded}{\begin{snugshade}}{\end{snugshade}}
\newcommand{\AlertTok}[1]{\textcolor[rgb]{0.68,0.00,0.00}{#1}}
\newcommand{\AnnotationTok}[1]{\textcolor[rgb]{0.37,0.37,0.37}{#1}}
\newcommand{\AttributeTok}[1]{\textcolor[rgb]{0.40,0.45,0.13}{#1}}
\newcommand{\BaseNTok}[1]{\textcolor[rgb]{0.68,0.00,0.00}{#1}}
\newcommand{\BuiltInTok}[1]{\textcolor[rgb]{0.00,0.23,0.31}{#1}}
\newcommand{\CharTok}[1]{\textcolor[rgb]{0.13,0.47,0.30}{#1}}
\newcommand{\CommentTok}[1]{\textcolor[rgb]{0.37,0.37,0.37}{#1}}
\newcommand{\CommentVarTok}[1]{\textcolor[rgb]{0.37,0.37,0.37}{\textit{#1}}}
\newcommand{\ConstantTok}[1]{\textcolor[rgb]{0.56,0.35,0.01}{#1}}
\newcommand{\ControlFlowTok}[1]{\textcolor[rgb]{0.00,0.23,0.31}{\textbf{#1}}}
\newcommand{\DataTypeTok}[1]{\textcolor[rgb]{0.68,0.00,0.00}{#1}}
\newcommand{\DecValTok}[1]{\textcolor[rgb]{0.68,0.00,0.00}{#1}}
\newcommand{\DocumentationTok}[1]{\textcolor[rgb]{0.37,0.37,0.37}{\textit{#1}}}
\newcommand{\ErrorTok}[1]{\textcolor[rgb]{0.68,0.00,0.00}{#1}}
\newcommand{\ExtensionTok}[1]{\textcolor[rgb]{0.00,0.23,0.31}{#1}}
\newcommand{\FloatTok}[1]{\textcolor[rgb]{0.68,0.00,0.00}{#1}}
\newcommand{\FunctionTok}[1]{\textcolor[rgb]{0.28,0.35,0.67}{#1}}
\newcommand{\ImportTok}[1]{\textcolor[rgb]{0.00,0.46,0.62}{#1}}
\newcommand{\InformationTok}[1]{\textcolor[rgb]{0.37,0.37,0.37}{#1}}
\newcommand{\KeywordTok}[1]{\textcolor[rgb]{0.00,0.23,0.31}{\textbf{#1}}}
\newcommand{\NormalTok}[1]{\textcolor[rgb]{0.00,0.23,0.31}{#1}}
\newcommand{\OperatorTok}[1]{\textcolor[rgb]{0.37,0.37,0.37}{#1}}
\newcommand{\OtherTok}[1]{\textcolor[rgb]{0.00,0.23,0.31}{#1}}
\newcommand{\PreprocessorTok}[1]{\textcolor[rgb]{0.68,0.00,0.00}{#1}}
\newcommand{\RegionMarkerTok}[1]{\textcolor[rgb]{0.00,0.23,0.31}{#1}}
\newcommand{\SpecialCharTok}[1]{\textcolor[rgb]{0.37,0.37,0.37}{#1}}
\newcommand{\SpecialStringTok}[1]{\textcolor[rgb]{0.13,0.47,0.30}{#1}}
\newcommand{\StringTok}[1]{\textcolor[rgb]{0.13,0.47,0.30}{#1}}
\newcommand{\VariableTok}[1]{\textcolor[rgb]{0.07,0.07,0.07}{#1}}
\newcommand{\VerbatimStringTok}[1]{\textcolor[rgb]{0.13,0.47,0.30}{#1}}
\newcommand{\WarningTok}[1]{\textcolor[rgb]{0.37,0.37,0.37}{\textit{#1}}}

\usepackage{longtable,booktabs,array}
\newcounter{none} % for unnumbered tables
\usepackage{calc} % for calculating minipage widths
\usepackage{caption}
% Make caption package work with longtable
\makeatletter
\def\fnum@table{\tablename~\thetable}
\makeatother
\usepackage{graphicx}
\makeatletter
\newsavebox\pandoc@box
\newcommand*\pandocbounded[1]{% scales image to fit in text height/width
  \sbox\pandoc@box{#1}%
  \Gscale@div\@tempa{\textheight}{\dimexpr\ht\pandoc@box+\dp\pandoc@box\relax}%
  \Gscale@div\@tempb{\linewidth}{\wd\pandoc@box}%
  \ifdim\@tempb\p@<\@tempa\p@\let\@tempa\@tempb\fi% select the smaller of both
  \ifdim\@tempa\p@<\p@\scalebox{\@tempa}{\usebox\pandoc@box}%
  \else\usebox{\pandoc@box}%
  \fi%
}
% Set default figure placement to htbp
\def\fps@figure{htbp}
\makeatother





\setlength{\emergencystretch}{3em} % prevent overfull lines

\providecommand{\tightlist}{%
  \setlength{\itemsep}{0pt}\setlength{\parskip}{0pt}}



 


\setbeamertemplate{navigation symbols}{}
\setbeamertemplate{footline}[frame number]
\setbeamertemplate{itemize items}[circle]
\setbeamertemplate{itemize subitem}[triangle]
\setbeamercolor{frametitle}{fg=black}
\setbeamerfont{frametitle}{series=\bfseries}
\usepackage{amsmath}
\newcommand{\indep}{\perp\!\!\!\perp}
\usepackage{booktabs}
\usepackage{array}
\newcolumntype{R}[1]{>{\raggedright\arraybackslash}p{#1}}
\usepackage{fancyvrb}
\fvset{fontsize=\scriptsize}
\makeatletter
\def\verbatim@font{\scriptsize\ttfamily}
\makeatother
\makeatletter
\@ifpackageloaded{caption}{}{\usepackage{caption}}
\AtBeginDocument{%
\ifdefined\contentsname
  \renewcommand*\contentsname{Table of contents}
\else
  \newcommand\contentsname{Table of contents}
\fi
\ifdefined\listfigurename
  \renewcommand*\listfigurename{List of Figures}
\else
  \newcommand\listfigurename{List of Figures}
\fi
\ifdefined\listtablename
  \renewcommand*\listtablename{List of Tables}
\else
  \newcommand\listtablename{List of Tables}
\fi
\ifdefined\figurename
  \renewcommand*\figurename{Figure}
\else
  \newcommand\figurename{Figure}
\fi
\ifdefined\tablename
  \renewcommand*\tablename{Table}
\else
  \newcommand\tablename{Table}
\fi
}
\@ifpackageloaded{float}{}{\usepackage{float}}
\floatstyle{ruled}
\@ifundefined{c@chapter}{\newfloat{codelisting}{h}{lop}}{\newfloat{codelisting}{h}{lop}[chapter]}
\floatname{codelisting}{Listing}
\newcommand*\listoflistings{\listof{codelisting}{List of Listings}}
\makeatother
\makeatletter
\makeatother
\makeatletter
\@ifpackageloaded{caption}{}{\usepackage{caption}}
\@ifpackageloaded{subcaption}{}{\usepackage{subcaption}}
\makeatother

\usepackage{bookmark}
\IfFileExists{xurl.sty}{\usepackage{xurl}}{} % add URL line breaks if available
\urlstyle{same}
\hypersetup{
  pdftitle={Synthetic Control},
  pdfauthor={Professor Ji-Woong Chung},
  hidelinks,
  pdfcreator={LaTeX via pandoc}}


\title{\texorpdfstring{Synthetic Control}{Synthetic Control}}
\subtitle{\texorpdfstring{BUSS975: Causal Inference in Financial
Research}{BUSS975: Causal Inference in Financial Research}}
\author{Professor Ji-Woong Chung}
\date{}
\institute{Korea University}

\begin{document}
\frame{\titlepage}


\begin{frame}{Where we left off}
\protect\phantomsection\label{where-we-left-off}
\begin{itemize}
\tightlist
\item
  Lectures 7a--7b: difference-in-differences identifies the ATT if
  untreated outcomes move in parallel. The staggered-adoption estimators
  refuse the forbidden comparison, but every one of them \textbf{still
  needs parallel trends for a control group you chose}.
\item
  Two questions stayed open:

  \begin{itemize}
  \tightlist
  \item
    What if the treated ``group'' is \textbf{one unit}: a state, a
    country, an exchange? One unit cannot choose its control group by
    regression, and there is nothing to cluster on.
  \item
    What if no untreated group is parallel to it, but \textbf{some
    weighted average of them} might be?
  \end{itemize}
\end{itemize}

\vspace{0.4em}

\begin{itemize}
\tightlist
\item
  Synthetic control (Abadie and Gardeazabal 2003; Abadie, Diamond and
  Hainmueller 2010) answers both. It imputes the treated unit's
  untreated outcome from a \textbf{data-chosen, non-negative weighted
  average of untreated units}. It targets the ATT and draws an
  event-study picture, like DiD. It does \textbf{not} assume parallel
  trends; it assumes a model of the outcome.
\item
  Source: Cunningham (2025), \emph{The Remix}, ch.~11, and his two
  datasets: Proposition 99 (ADH 2010) and Texas prisons. Every estimate
  here is computed by hand.
\end{itemize}
\end{frame}

\begin{frame}{The plan}
\protect\phantomsection\label{the-plan}
\begin{enumerate}
\tightlist
\item
  \textbf{The comparative case study.} Card's Mariel boatlift and its
  two weaknesses.
\item
  \textbf{The estimator.} Weights, the \(V\) matrix, three advantages,
  the factor model.
\item
  \textbf{Proposition 99.} Synthetic California, balance, and inference
  by placebo.
\item
  \textbf{Texas prisons.} The data exercise, placebos in space and in
  time, and two lessons.
\item
  \textbf{Weights, extrapolation, imperfect fit.} Regression weights,
  the convex hull, specification choice, demeaning, augmentation.
\item
  \textbf{Matrix completion.} The potential-outcome matrix, low rank,
  and the nuclear norm.
\item
  \textbf{Synthetic difference-in-differences.} Unit weights, time
  weights, and a weighted DiD.
\item
  \textbf{Using it} in finance.
\end{enumerate}

\vspace{0.4em}

\begin{itemize}
\tightlist
\item
  One running question: \textbf{which untreated units are inside my
  counterfactual, with what weight, and what do I give up to fit
  better?}
\end{itemize}
\end{frame}

\section{1. The Comparative Case
Study}\label{the-comparative-case-study}

\begin{frame}{Card (1990): the Mariel boatlift}
\protect\phantomsection\label{card-1990-the-mariel-boatlift}
\begin{itemize}
\tightlist
\item
  April--October 1980: Castro opens the port of Mariel; about 125,000
  Cubans leave, most settle in Miami. Miami's labour force grows by
  about 7\%, mostly low-skilled workers. A textbook supply shock in a
  concentrated local labour market.
\item
  Card compares Miami with four cities (Atlanta, Houston, Los Angeles,
  Tampa--St.~Petersburg) in a difference-in-differences on CPS wages and
  unemployment, and finds \textbf{no discernible effect}. The four
  cities are chosen in a footnote for ``similar demographics and
  economic conditions.''
\item
  Controversial because of the answer, not the method. Peri and Yasenov
  (2019) redo it with synthetic control and reach the same conclusion.
\end{itemize}

\vspace{0.4em}

\begin{itemize}
\tightlist
\item
  Two weaknesses of the \textbf{design}, whatever one thinks of the
  finding:

  \begin{enumerate}
  \tightlist
  \item
    The control group was \textbf{ad hoc}. A different set of cities
    might give a different answer, and the reader cannot tell how much
    the choice mattered.
  \item
    The standard errors measured \textbf{sampling variance}, not the
    uncertainty that matters: whether those four cities reproduce what
    Miami would have done without the boatlift.
  \end{enumerate}
\item
  Synthetic control was built to address both at once.
\end{itemize}
\end{frame}

\begin{frame}{What the method does}
\protect\phantomsection\label{what-the-method-does}
\begin{itemize}
\tightlist
\item
  One treated aggregate unit, \(j = 1\), and a \textbf{donor pool} of
  \(J\) untreated units, \(j = 2, \dots, J+1\), observed for \(T_0\)
  pre-treatment periods and some post periods.
\item
  The synthetic control is a weighted average of donors with
  \(w_j \geq 0\) and \(\sum_j w_j = 1\), the weights chosen so that it
  \textbf{resembles the treated unit before treatment}. The effect in
  post period \(t\) is \[
  \hat\tau_{1t} = Y_{1t} - \sum_{j=2}^{J+1} w_j^{*}\, Y_{jt}.
  \]
\item
  Target: the ATT for the treated unit, period by period. The estimator
  is a picture: two series that coincide before \(T_0\) and separate
  after it, or do not.
\item
  Shared with DiD: panel data, the ATT, an event-study plot, extensions
  to several treated units and staggered timing. Not shared:
  \textbf{parallel trends}. The comparison group is chosen by fit, not
  by assumption, and Cunningham keeps it outside the DiD family tree.
\item
  Lineage: Abadie was one of Angrist's first students and a long-time
  coauthor of Imbens. Like the LATE, the method sits where Princeton
  labour economics and Harvard statistics meet.
\end{itemize}
\end{frame}

\section{2. The Estimator}\label{the-estimator}

\begin{frame}{Choosing the weights}
\protect\phantomsection\label{choosing-the-weights}
\begin{itemize}
\tightlist
\item
  \(X_1\) (\(k \times 1\)): pre-treatment characteristics of the treated
  unit; \(X_0\) (\(k \times J\)): the same for donors. Characteristics
  are \textbf{predictors of post-treatment outcomes}: covariates and
  lagged outcomes, chosen by the researcher.
\item
  ADH (2010) choose \(W = (w_2, \dots, w_{J+1})'\) to minimise \[
  \|X_1 - X_0 W\|_V = \sqrt{(X_1 - X_0 W)'\, V\, (X_1 - X_0 W)}
  \quad\text{s.t.}\quad w_j \geq 0,\ \ \textstyle\sum_j w_j = 1 .
  \]
\item
  With diagonal \(V\) the objective is a \textbf{matching discrepancy},
  \(\sum_m v_m \big(X_{1m} - \sum_j w_j X_{jm}\big)^2\): Lecture 3's
  nearest-neighbour matching, with one treated unit matched to a convex
  combination of controls.
\item
  \(V\) is the \textbf{importance matrix}, and \(W^{*}(V)\) depends on
  it. The usual choice: \(V\) that minimises the pre-period root mean
  squared prediction error of the outcome, \[
  \text{RMSPE} = \Big( \tfrac{1}{T_0} \textstyle\sum_{t=1}^{T_0} \big( Y_{1t} - \sum_j w_j^{*}(V)\, Y_{jt} \big)^2 \Big)^{1/2},
  \] a \textbf{nested} optimisation: a quadratic programme for \(W\)
  inside a search over \(V\).
\item
  ``Optimal'' means no other non-negative, sum-to-one vector matches
  \(X_1\) better. Drop the constraints and you can always match better;
  Section 5 is about that trade.
\end{itemize}
\end{frame}

\begin{frame}{Three advantages (Abadie, Diamond and Hainmueller 2010)}
\protect\phantomsection\label{three-advantages-abadie-diamond-and-hainmueller-2010}
\begin{enumerate}
\tightlist
\item
  \textbf{No extrapolation.} The counterfactual is a convex combination
  of \emph{contemporaneous} donor outcomes. It never projects forward
  from the treated unit's own history and never steps outside the data,
  where regression is most model-dependent (King and Zeng 2006).
\item
  \textbf{Design before outcomes.} The weights use only pre-treatment
  data. Fix the model, publish the weights, and only then look at the
  post period (Rubin 2007, 2008): the ``do not peek'' principle of the
  7b checklist, built into the estimator.
\item
  \textbf{Transparency.} The weights say exactly what each donor
  contributes. A regression also weights the donors, but the weights are
  buried in \((X'X)^{-1}X'\) and nobody shows them.
\end{enumerate}

\vspace{0.4em}

\begin{itemize}
\tightlist
\item
  The seminar objection: ``Idaho gets 0.3 of California? Absurd!'' In a
  regression Idaho might get a \emph{larger} weight, or a negative one,
  and no one would know (ADH 2015). Transparency invites scrutiny that
  other estimators escape. That is a strength.
\item
  The cost of (1), seen later: when no convex combination resembles the
  treated unit, the method cannot fit, and it says so.
\end{itemize}
\end{frame}

\begin{frame}{The outcome model}
\protect\phantomsection\label{the-outcome-model}
\begin{itemize}
\tightlist
\item
  The assumption behind the method is a \textbf{factor model} for
  untreated outcomes: \[
  Y^0_{it} = \alpha_t + \theta_t Z_i + \lambda_t u_i + \varepsilon_{it},
  \] \(\alpha_t\) a common shock; \(Z_i\) observed characteristics with
  \textbf{time-varying} coefficients \(\theta_t\); \(u_i\) unobserved
  loadings on unobserved common factors \(\lambda_t\);
  \(\varepsilon_{it}\) transitory noise.
\item
  A model of the \textbf{outcome}, not of assignment. Propensity scores
  model who gets treated; synthetic control, like the outcome regression
  of 7b, models what the untreated outcome would have been, with a
  stronger structure: a few latent factors. It is 7b's
  ``covariate-specific trend'' in other clothes: grit is stable, the
  return to grit moves. In DiD that breaks parallel trends; here it is
  the model.
\item
  If \(\sum_j w_j Z_j = Z_1\) \textbf{and} \(\sum_j w_j u_j = u_1\),
  then \(\sum_j w_j Y^0_{jt} = Y^0_{1t}\) up to noise in every period,
  treated or not. \(u_i\) is unobserved, so the \textbf{pre-period fit
  is the evidence}: only units balanced on both \(Z\) and \(u\) would
  track the treated unit for many years. ADH show the bias vanishes as
  \(T_0 \to \infty\) when the fit is exact.
\item
  Hence the rule: a \textbf{long} pre-period, \textbf{closely} tracked.
  The assumption is not testable; the fit is imperfect but powerful
  evidence.
\end{itemize}
\end{frame}

\section{3. Proposition 99}\label{proposition-99}

\begin{frame}{California, 1988}
\protect\phantomsection\label{california-1988}
\begin{center}
\includegraphics[width=0.78\linewidth,height=\textheight,keepaspectratio]{08_synthetic_control_files/figure-beamer/prop99-raw-1.pdf}
\end{center}

\begin{itemize}
\tightlist
\item
  Proposition 99 (passed November 1988, effective January 1989): a
  25-cent tax per pack, clean-air ordinances, anti-smoking media
  campaigns, more than \$100 million a year earmarked for anti-tobacco
  programmes.
\item
  Outcome: cigarette sales per capita, 1970--2000: 19 pre-years, 12
  post-years. Donor pool: the 38 states without comparable programmes in
  the period; states with similar legislation are dropped, for the
  reason 7b gave against comparing with the already-treated.
\item
  Sales fell after 1988. They were already falling, in California and
  everywhere else. The raw comparison cannot tell.
\end{itemize}
\end{frame}

\begin{frame}{Synthetic California}
\protect\phantomsection\label{synthetic-california}
\begin{center}
\includegraphics[width=0.78\linewidth,height=\textheight,keepaspectratio]{08_synthetic_control_files/figure-beamer/prop99-synth-1.pdf}
\end{center}

\begin{itemize}
\tightlist
\item
  Predictors (ADH 2010): ln GDP per capita, share aged 15--24, retail
  price, beer consumption (1980--88 averages; beer 1984--88), and
  cigarette sales in 1975, 1980 and 1988.
\item
  Weights: Colorado 0.136, Connecticut 0.059, Montana 0.219, Nevada
  0.242, Utah 0.344; the other 33 states get zero. ADH report Utah
  0.334, Nevada 0.234, Montana 0.199, Colorado 0.164, Connecticut 0.069.
\item
  Pre-period RMSPE 1.78 packs. \(V\) concentrates on retail price
  (0.71), cigarette sales 1980 (0.21), cigarette sales 1975 (0.05); the
  other predictors get almost nothing.
\end{itemize}
\end{frame}

\begin{frame}{Balance}
\protect\phantomsection\label{balance}
\begin{center}\small
\begin{tabular}{lrrr}
\toprule
 & California & Synthetic & Average of 38 \\
\midrule
Ln(GDP per capita)  &  10.08  &  9.85  &  9.83  \\
Percent aged 15--24  &  17.35  &  17.35  &  17.25  \\
Retail price  &  89.42  &  89.42  &  87.27  \\
Beer consumption per capita  &  24.28  &  24.27  &  23.66  \\
Cigarette sales 1988  &  90.10  &  91.36  &  113.82  \\
Cigarette sales 1980  &  120.20  &  120.20  &  138.09  \\
Cigarette sales 1975  &  127.10  &  127.09  &  136.93  \\
\bottomrule
\end{tabular}
\end{center}

\begin{itemize}
\tightlist
\item
  One number per cell, so this is not a test; it is the best the method
  allows. The synthetic control matches lagged sales to within two
  packs; the donor average misses 1988 sales by 23.7.
\item
  Income is the one predictor that is not matched: \(V\) gave it no
  weight, because it did not help predict pre-period sales once the lags
  were in. That is what \(V\) is for.
\end{itemize}
\end{frame}

\begin{frame}{The gap}
\protect\phantomsection\label{the-gap}
\begin{center}
\includegraphics[width=0.78\linewidth,height=\textheight,keepaspectratio]{08_synthetic_control_files/figure-beamer/prop99-gap-1.pdf}
\end{center}

\begin{itemize}
\tightlist
\item
  The estimator is the difference between the two series: near zero for
  19 years, then a widening gap. Mean over 1989--2000: \(-18.9\) packs
  per capita; by 2000, \(-25.6\) packs, about a quarter of what
  California would have sold.
\item
  Like RD, this is a picture-intensive estimator. If the picture does
  not show flat-then-diverge, no table will rescue it.
\end{itemize}
\end{frame}

\begin{frame}{Inference: ``no treatment effect whatsoever''}
\protect\phantomsection\label{inference-no-treatment-effect-whatsoever}
\begin{itemize}
\tightlist
\item
  Two observations per year and one treated unit: nothing to cluster on.
  Maybe the divergence is prediction error that any model would produce.
\item
  ADH (2010): \textbf{randomization inference} (Fisher 1935) under the
  sharp null of no effect for any unit in any period, ``no treatment
  effect whatsoever'' (Firpo and Possebom 2018). Lecture 1's idea, with
  the donor pool as the randomization distribution.
\item
  The test statistic is the ratio of post- to pre-treatment RMSPE,
  \(r_j = \text{RMSPE}_{\text{post}, j} / \text{RMSPE}_{\text{pre}, j}\),
  with \[
  \text{RMSPE}_{\text{post}} = \Big( \tfrac{1}{T - T_0} \textstyle\sum_{t = T_0 + 1}^{T} \big( Y_{1t} - \sum_j w_j^{*} Y_{jt} \big)^2 \Big)^{1/2}.
  \]
\item
  Steps: (1) treat each donor in turn as if it had been treated, with
  California back in the pool, and refit; (2) pre RMSPE; (3) post RMSPE;
  (4) the ratio; (5) sort; (6) \(p = \text{rank}/\text{total}\).
\item
  Why a ratio: a state the method could not fit before treatment shows a
  large post gap for free. Dividing by its own pre-period fit asks
  whether the post gap is \textbf{unusual relative to the fit}.
\end{itemize}
\end{frame}

\begin{frame}{Placebo gaps: all 38 donors}
\protect\phantomsection\label{placebo-gaps-all-38-donors}
\begin{center}
\includegraphics[width=0.78\linewidth,height=\textheight,keepaspectratio]{08_synthetic_control_files/figure-beamer/prop99-placebo-all-1.pdf}
\end{center}

\begin{itemize}
\tightlist
\item
  California is the bold line; each grey line is a donor state
  pretending to be treated in 1989, with its own synthetic control.
\item
  California's post gap is among the largest, but a few placebos have
  wild paths before and after 1989. Those are states the method cannot
  fit: extreme units with no convex combination of neighbours. They blow
  up the scale.
\end{itemize}
\end{frame}

\begin{frame}{Placebo gaps: states fit as well as California}
\protect\phantomsection\label{placebo-gaps-states-fit-as-well-as-california}
\begin{center}
\includegraphics[width=0.78\linewidth,height=\textheight,keepaspectratio]{08_synthetic_control_files/figure-beamer/prop99-placebo-good-1.pdf}
\end{center}

\begin{itemize}
\tightlist
\item
  ADH drop placebos whose pre-period RMSPE exceeds some multiple of
  California's; here twice, which keeps 22 of 38 donors. Among units the
  method fits as well as it fits California, none opens a gap of
  California's size; the largest placebo gap is about half of it.
\item
  A judgement call, so show the cut-offs you tried. The ratio test on
  the next slide needs no cut-off.
\end{itemize}
\end{frame}

\begin{frame}[fragile]{The exact p-value}
\protect\phantomsection\label{the-exact-p-value}
\begin{center}
\includegraphics[width=0.72\linewidth,height=\textheight,keepaspectratio]{08_synthetic_control_files/figure-beamer/prop99-hist-1.pdf}
\end{center}

\begin{itemize}
\tightlist
\item
  California's ratio is 11.2. ADH, with the \texttt{Synth} optimiser,
  rank it first of 39: \(p = 1/39 = 0.026\). This deck's optimiser ranks
  it 2: \(p = 2/39 = 0.051\). The state above it is Missouri (ratio
  12.6), which the method fits even more closely before 1989 (pre-period
  RMSPE 0.83) and which drifts 9.7 packs from its synthetic control
  after.
\item
  The lesson is not that the test failed. The ratio divides by a
  denominator that the \(V\) search chose, and a slightly different
  \(V\) gives one placebo a tighter fit and a larger ratio. Report the
  histogram, the treated unit's rank, and the optimiser; and do not read
  0.026 against 0.051 as a difference in kind.
\end{itemize}
\end{frame}

\section{4. Texas Prisons}\label{texas-prisons}

\begin{frame}{Ruiz v. Estelle and the construction boom}
\protect\phantomsection\label{ruiz-v.-estelle-and-the-construction-boom}
\begin{center}
\includegraphics[width=0.72\linewidth,height=\textheight,keepaspectratio]{08_synthetic_control_files/figure-beamer/texas-capacity-1.pdf}
\end{center}

\begin{itemize}
\tightlist
\item
  1980: inmate David Ruiz wins a suit against the Texas Department of
  Corrections over overcrowding. Courts cap inmates per cell; federal
  supervision lasts until 2003. Through the 1980s Texas manages
  admissions with \textbf{early parole}, since it can neither overcrowd
  nor afford to build.
\item
  Modest construction under Governor Clements in the late 1980s. In
  1993, under Governor Richards, a billion-dollar package: capacity
  grows by more than 30\% a year for three years and \textbf{doubles}.
\item
  Capacity is not incarceration. But parole had been the relief valve;
  once the beds existed, Texas kept people it would have released. The
  group affected most: Black men.
\end{itemize}
\end{frame}

\begin{frame}[fragile]{Texas against the rest of the country}
\protect\phantomsection\label{texas-against-the-rest-of-the-country}
\begin{center}
\includegraphics[width=0.78\linewidth,height=\textheight,keepaspectratio]{08_synthetic_control_files/figure-beamer/texas-raw-1.pdf}
\end{center}

\begin{itemize}
\tightlist
\item
  Black male prisoners in Texas roughly double in two years after 1993.
  The rest of the country rises smoothly.
\item
  Before 1993 Texas was \textbf{unremarkable}: run-of-the-mill rather
  than the largest or smallest state on this outcome. That is what makes
  a non-negative weighted average of other states feasible. Had Texas
  been extreme, no convex combination could have matched it.
\item
  Cunningham uses the count, not the rate; the rate is in the data and
  is the chapter's exercise. The analysis follows his
  \texttt{synth\_1.do}.
\end{itemize}
\end{frame}

\begin{frame}{Synthetic Texas}
\protect\phantomsection\label{synthetic-texas}
\begin{center}
\includegraphics[width=0.78\linewidth,height=\textheight,keepaspectratio]{08_synthetic_control_files/figure-beamer/texas-synth-1.pdf}
\end{center}

\begin{itemize}
\tightlist
\item
  Predictors: Black male prisoners in 1988, 1990, 1991, 1992; alcohol
  consumption 1990; AIDS cases per capita 1990, 1991; income,
  unemployment and poverty (pre-period averages); percent Black
  1990--92; percent aged 15--19 in 1990. Fit period 1985--1993, as in
  the chapter.
\item
  Weights: Florida 0.457, Illinois 0.235, New York 0.308. Pre-period
  RMSPE 842 prisoners against a 1993 level of 29,260.
\item
  This is the picture one hopes for: a synthetic Texas on top of Texas
  for nine years, then a divergence as the beds open.
\end{itemize}
\end{frame}

\begin{frame}{The gap and the balance}
\protect\phantomsection\label{the-gap-and-the-balance}
\begin{center}
\includegraphics[width=0.72\linewidth,height=\textheight,keepaspectratio]{08_synthetic_control_files/figure-beamer/texas-gap-1.pdf}
\end{center}

\begin{center}\scriptsize
\begin{tabular}{lrrr}
\toprule
 & Texas & Synthetic & Average of 50 \\
\midrule
Black male prisoners 1990  &  22,634  &  22,720  &  6,098  \\
Black male prisoners 1992  &  27,568  &  25,861  &  6,972  \\
Alcohol 1990  &  2  &  3  &  3  \\
Income  &  16,486  &  19,534  &  17,480  \\
Unemployment  &  7  &  7  &  6  \\
Percent Black 1990  &  16  &  16  &  11  \\
Percent aged 15--19, 1990  &  8  &  7  &  7  \\
\bottomrule
\end{tabular}
\end{center}

\begin{itemize}
\tightlist
\item
  Mean gap over 1994--2000: 23,612 prisoners; in 2000, 28,466. Lagged
  counts, alcohol and the Black population share are matched; income and
  the age share are not.
\end{itemize}
\end{frame}

\begin{frame}{Placebos in space}
\protect\phantomsection\label{placebos-in-space}
\begin{center}
\includegraphics[width=0.78\linewidth,height=\textheight,keepaspectratio]{08_synthetic_control_files/figure-beamer/texas-placebo-1.pdf}
\end{center}

\begin{itemize}
\tightlist
\item
  Left: gaps for every state with a pre-period RMSPE within twice that
  of Texas (49 of 50); Texas in bold. Right: the ratio for all 51.
\item
  Texas's ratio is 28.8, rank 1 of 51: \(p = 0.020\). The runner-up is
  Wisconsin at 20.9.
\item
  Same recipe as California: refit 51 times, pretend each state built
  the prisons, compare ratios.
\end{itemize}
\end{frame}

\begin{frame}{A placebo in time}
\protect\phantomsection\label{a-placebo-in-time}
\begin{center}
\includegraphics[width=0.78\linewidth,height=\textheight,keepaspectratio]{08_synthetic_control_files/figure-beamer/texas-intime-1.pdf}
\end{center}

\begin{itemize}
\tightlist
\item
  ADH (2015) recommend a second falsification: move the treatment date
  to a year when nothing happened and stop the sample before the real
  treatment. Here 1989, with 1985--1992 and predictors shifted back
  accordingly.
\item
  A gap opening after 1989 would mean the method manufactures effects
  from this donor pool. The 1990--92 gap averages 1,486 prisoners,
  against 23,612 after the real date.
\item
  The two placebos answer different questions: in space, ``is Texas
  unusual among states?''; in time, ``is 1993 unusual among years?''
\end{itemize}
\end{frame}

\begin{frame}[fragile]{Two lessons from the exercise}
\protect\phantomsection\label{two-lessons-from-the-exercise}
\begin{itemize}
\tightlist
\item
  \textbf{The \(V\) step is not convex.} The inner problem is a
  quadratic programme with a unique answer; the search over \(V\) has
  many local minima. This deck's optimiser lands on Florida, Illinois,
  New York; Stata's \texttt{synth} on the same specification lands on
  California 0.41, Illinois 0.36, Louisiana 0.12, Florida 0.11, with a
  similar fit and gap. Report the fit and the gap, not only the weights;
  try several starting values.
\item
  \textbf{Specification matters.} Which lags, which covariates, which
  fit period. Ferman, Pinto and Possebom (2020): searching over these
  raises false positives, since some specification will usually produce
  a gap. Their remedy, for Proposition 99:
\end{itemize}

\begin{center}\scriptsize
\begin{tabular}{lrr}
\toprule
Predictors & Pre RMSPE & Mean gap 1989--2000 \\
\midrule
ADH: four covariates + sales in 1975, 1980, 1988  &  1.78  &  -18.9  \\
All 19 pre-period sales  &  1.66  &  -19.5  \\
Sales in 1975, 1980, 1988 only  &  3.99  &  -21.2  \\
Odd years only  &  1.87  &  -20.3  \\
Four covariates + mean pre-period sales  &  4.71  &  -20.0  \\
Demeaned (with intercept)  &  0.96  &  -11.1  \\
Regression weights  &  2.08  &  -11.2  \\
\bottomrule
\end{tabular}
\end{center}
\end{frame}

\section{5. Weights, Extrapolation, and Imperfect
Fit}\label{weights-extrapolation-and-imperfect-fit}

\begin{frame}{Synthetic weights against regression weights}
\protect\phantomsection\label{synthetic-weights-against-regression-weights}
\begin{itemize}
\tightlist
\item
  Regress post-treatment outcomes on the \(k\) predictors and a
  constant, among donors, and predict for the treated unit. That too
  imputes \(Y^0_{1t}\) as a weighted average of donors (Abadie, Diamond
  and Hainmueller 2015): \[
  w^{\text{reg}} = \tilde X_0' \big( \tilde X_0 \tilde X_0' \big)^{-1} \tilde X_1, \qquad \tilde X = \begin{pmatrix} 1 \\ X \end{pmatrix}.
  \] The weights sum to one because of the constant. They need not be
  non-negative, and they are rarely zero.
\item
  Proposition 99 with the seven ADH predictors: 36 of 38 states carry
  weight above 0.01, 12 are negative, the smallest is Arkansas at
  \(-0.14\). Synthetic control uses 5 states, all positive.
\item
  ADH (2015), West German reunification, 16 donors: synthetic control
  uses Austria 0.42, United States 0.22, Japan 0.16, Switzerland 0.11,
  Netherlands 0.09; regression weights every country, Greece at
  \(-0.09\) and Portugal at \(-0.08\).
\item
  Regression has more freedom, so it fits at least as well; it buys the
  fit with \textbf{extrapolation} along an assumed functional form,
  outside the support of the donors. Here the regression counterfactual
  gives a mean gap of \(-11.2\) packs.
\end{itemize}
\end{frame}

\begin{frame}{Interpolation and the convex hull}
\protect\phantomsection\label{interpolation-and-the-convex-hull}
\begin{itemize}
\tightlist
\item
  \(w_j \geq 0\) and \(\sum w_j = 1\) mean the synthetic control lies in
  the \textbf{convex hull} of the donors. Synthetic control is matching,
  and the convex hull is Lecture 3's common support in panel clothes.
\item
  No convex combination of basketball players is Michael Jordan; none of
  footballers is Messi. For an outlier, the ``optimal'' non-negative
  weights are the best available under the constraints, and the fit is
  poor. Only \textbf{negative weights}, anti-Jordans, could close the
  gap, and that is extrapolation.
\item
  The limitation is the strength. By refusing extreme counterfactuals,
  synthetic control gives up some fit and escapes model dependence (King
  and Zeng 2006).
\end{itemize}

\vspace{0.4em}

\begin{itemize}
\tightlist
\item
  When the fit is imperfect, synthetic control is \textbf{biased}: the
  weights are tracking transitory noise \(\varepsilon\) rather than the
  loadings \(u\). The bias bound grows with the noise variance and
  shrinks with \(T_0\) (ADH 2010; Ferman and Pinto 2021). ADH (2015):
  ``We do not recommend using this method when the pretreatment fit is
  poor or the number of pretreatment periods is small.''
\item
  This is not an argument for long pre-periods in DiD, which identifies
  the ATT with one pre-period and needs more only for credibility.
\end{itemize}
\end{frame}

\begin{frame}{Covariates, and what they are for}
\protect\phantomsection\label{covariates-and-what-they-are-for}
\begin{itemize}
\tightlist
\item
  In unconfoundedness, covariates close back doors; in DiD, they make
  parallel trends conditional. In synthetic control their only job is to
  \textbf{align donors with the treated unit before treatment}, so that
  the weights recover the factor loadings. They are worth including
  insofar as they do that.
\item
  Lagged outcomes already carry \(\theta_t Z_i + \lambda_t u_i\).
  Increasingly researchers use \textbf{only lagged outcomes}
  (Ben-Michael, Feller and Rothstein 2021), as the second-generation
  estimators in this lecture do. ADH's covariate-plus-three-lags
  specification and the all-lags specification gave \(-18.9\) and
  \(-19.5\) for Proposition 99.
\item
  Kaul, Klößner, Pfeifer and Schieler (2022): if \textbf{all} pre-period
  outcomes are included as predictors, \(V\) puts zero weight on every
  covariate, and the covariates are irrelevant. Including all lags and a
  long covariate list is decoration.
\item
  The choice is a specification choice and should be reported as one.
  The table two slides back is the format: predictors in rows, fit and
  estimate in columns. If the estimate survives every sensible column,
  the covariate question is settled; if not, the fit column tells you
  which specifications to trust.
\end{itemize}
\end{frame}

\begin{frame}{Demeaned synthetic control (Ferman and Pinto 2021)}
\protect\phantomsection\label{demeaned-synthetic-control-ferman-and-pinto-2021}
\begin{itemize}
\tightlist
\item
  Subtract each unit's own pre-period mean, fit the weights to the
  \textbf{demeaned} series, add the treated unit's mean back: \[
  \hat Y^0_{1t} = \bar Y_{1,\text{pre}} + \sum_j w_j^{\text{dm}} \big( Y_{jt} - \bar Y_{j,\text{pre}} \big).
  \] This is synthetic control \textbf{with an intercept}, the first
  relaxation proposed by Doudchenko and Imbens (2017).
\item
  What it does: removes the level, as unit fixed effects do in DiD. The
  treated unit need no longer lie inside the convex hull in levels, only
  in deviations. Shared exposure to the aggregate factors \(\lambda_t\)
  is matched on the dynamics, and the level difference is absorbed,
  exactly as a DiD absorbs a level gap.
\item
  Proposition 99: pre-period RMSPE falls from 1.66 to 0.96; the weights
  spread over 9 states instead of 6; the estimate moves from \(-19.5\)
  to \(-11.1\) packs. The fit improved and the answer changed: the level
  constraint had been binding.
\item
  Less bias under imperfect fit, more variance in small samples, and a
  different assumption, closer to DiD's. Ferman and Pinto give a
  specification test for when it is appropriate.
\end{itemize}
\end{frame}

\begin{frame}{Augmented synthetic control}
\protect\phantomsection\label{augmented-synthetic-control}
\begin{itemize}
\tightlist
\item
  Ben-Michael, Feller and Rothstein (2021) apply the Abadie--Imbens
  (2011) correction for inexact matching: estimate the outcome model,
  subtract the imbalance. Write
  \(Y^0_{jt} = m_{jt} + \varepsilon_{jt}\), estimate \(\hat m(\cdot)\)
  from donors, and
\end{itemize}

\small

\[
\begin{aligned}
\hat Y^{\text{aug},0}_{1} &= \textstyle\sum_{D_j = 0} \hat w_j^{\text{synth}} Y_j + \underbrace{\big[ \hat m(X_1) - \sum_{D_j = 0} \hat w_j^{\text{synth}}\, \hat m(X_j) \big]}_{\text{estimated matching bias}}
= \hat m(X_1) + \sum_{D_j = 0} \hat w_j^{\text{synth}} \big( Y_j - \hat m(X_j) \big).
\end{aligned}
\] \normalsize

\begin{itemize}
\tightlist
\item
  The bracket is the outcome model's prediction of how far the synthetic
  control's imbalance in \(X\) translates into imbalance in \(Y^0\).
  With a perfect fit it is zero and augmented equals classic. With an
  imperfect fit it corrects in proportion to the imbalance.
\item
  Second form: the outcome regression, debiased by a weighted average of
  donor residuals. The doubly robust structure of 3b and 7b again.
\item
  Why bother, if \(T_0 \to \infty\) kills the bias? Because \(T_0\) is
  19 for smoking and 9 for Texas. Exact balance is rare.
\end{itemize}
\end{frame}

\begin{frame}{The ridge weights}
\protect\phantomsection\label{the-ridge-weights}
\begin{itemize}
\tightlist
\item
  \(\hat m\) is a prediction problem with \(T_0\) regressors (the lagged
  outcomes) and \(J\) observations, so it is regularised. Ridge, on
  donors:
  \(\min_{\eta_0, \eta}\ \tfrac{1}{2} \sum_{D_j = 0} ( Y_j - \eta_0 - X_j' \eta )^2 + \lambda \|\eta\|_2^2\).
\item
  Substituting \(\hat m(X) = \hat\eta_0 + X'\hat\eta\) gives a
  \textbf{single set of weights}:
\end{itemize}

\small

\[
\begin{aligned}
\hat Y^{\text{aug},0}_1 &= \textstyle\sum_{D_j = 0} \hat w_j^{\text{synth}} Y_j + \big( X_1 - \sum_{D_j = 0} \hat w_j^{\text{synth}} X_j \big)' \hat\eta = \sum_{D_j = 0} \hat w_j^{\text{aug}} Y_j, \\
\hat w^{\text{aug}} &= \hat w^{\text{synth}} + \big( X_1 - X_0' \hat w^{\text{synth}} \big)' \big( X_0' X_0 + \lambda I \big)^{-1} X_0' .
\end{aligned}
\] \normalsize

\begin{itemize}
\tightlist
\item
  The augmented weights can be negative. They are the \textbf{minimal}
  negative reweighting that moves the imbalance toward zero: large
  \(\lambda\) means little extrapolation and weights close to the
  synthetic ones. \(\lambda\) is chosen by cross-validation.
\item
  The bill: negative weights are extrapolation, and extrapolation is
  model dependence. The larger the imbalance, the larger the correction,
  and the more the estimate hinges on \(\hat m\). The more
  model-dependent you are, the further you are from a design.
\end{itemize}
\end{frame}

\begin{frame}[fragile]{Smoking and prisons, classic and augmented}
\protect\phantomsection\label{smoking-and-prisons-classic-and-augmented}
\begin{center}
\includegraphics[width=0.88\linewidth,height=\textheight,keepaspectratio]{08_synthetic_control_files/figure-beamer/ascm-1.pdf}
\end{center}

\begin{itemize}
\tightlist
\item
  Predictors: all pre-period outcomes, as in the chapter's
  \texttt{allsynth} runs. \textbf{Smoking:} pre-period RMSPE 1.66
  classic, 1.05 augmented; the mean gap moves from \(-19.5\) to
  \(-17.2\). Small but visible, so the augmented weights are not the
  classic ones: 18 of 38 are negative, the smallest \(-0.019\), the
  largest 0.39 (right panel, hollow = negative).
\item
  \textbf{Texas:} incarceration grew smoothly everywhere and Texas sat
  in the middle of the pack, so the classic fit is already good; the
  augmentation changes the mean gap by \(-2.2\)\%. When the fit is
  excellent the bias term has nothing to correct.
\end{itemize}
\end{frame}

\section{6. Matrix Completion}\label{matrix-completion}

\begin{frame}{From the Netflix prize to the potential-outcome matrix}
\protect\phantomsection\label{from-the-netflix-prize-to-the-potential-outcome-matrix}
\begin{itemize}
\tightlist
\item
  2006: Netflix offers \$1 million to whoever beats its recommendation
  engine by 10\%. 2009: BellKor's Pragmatic Chaos wins with an ensemble
  built around \textbf{matrix factorisation}: the user-by-film ratings
  matrix is mostly empty, and it can be filled because a few latent
  factors (taste, genre) drive it.
\item
  Athey, Bayati, Doudchenko, Imbens and Khosravi (2021) carry the idea
  to causal panels. Write untreated potential outcomes as an
  \(N \times T\) matrix. Treatment makes entries disappear: the
  switching equation turns \(Y^0\) off and \(Y^1\) on. \[
  \mathbf Y^0 = \begin{pmatrix} Y^0_{11} & Y^0_{12} & \cdots & Y^0_{1T} \\ Y^0_{21} & Y^0_{22} & \cdots & Y^0_{2T} \\ \vdots & & \ddots & \vdots \\ Y^0_{N1} & Y^0_{N2} & \cdots & Y^0_{NT} \end{pmatrix}
  \qquad\longrightarrow\qquad
  \begin{pmatrix} Y^0_{11} & Y^0_{12} & \cdots & ? \\ Y^0_{21} & Y^0_{22} & \cdots & ? \\ \vdots & & \ddots & \vdots \\ Y^0_{N1} & Y^0_{N2} & \cdots & Y^0_{NT} \end{pmatrix}
  \]
\item
  Every estimator in this course imputes the question marks. Matrix
  completion does it explicitly, using patterns \textbf{across units and
  over time at once}, for one treated unit or many, with simultaneous or
  staggered adoption.
\end{itemize}
\end{frame}

\begin{frame}{Vertical, horizontal; thin, fat, square}
\protect\phantomsection\label{vertical-horizontal-thin-fat-square}
\begin{itemize}
\tightlist
\item
  Doudchenko and Imbens (2017), Ferman and Pinto (2021): synthetic
  control is a \textbf{vertical regression}. The treated unit's
  pre-period outcomes are regressed on the donors' outcomes in the same
  periods; the \emph{columns} (periods) are the observations; the
  constraints are \(w \geq 0\) and no intercept.
\item
  Unconfoundedness is a \textbf{horizontal regression}: the last period
  on its own lags, within unit; the \emph{rows} are the observations.
  Its nonparametric version is matching on pre-treatment trajectories.
\item
  The shape of the matrix says which is natural:

  \begin{itemize}
  \tightlist
  \item
    \textbf{Thin} (\(N \gg T\)): LaLonde, thousands of units and one
    pre-year; a single treated period in the last column. Horizontal.
  \item
    \textbf{Fat} (\(T \gg N\)): synthetic control, 39 states and 31
    years; the treated row is missing its tail. Vertical.
  \item
    \textbf{Square, staggered}: treatment switches on at different
    dates, the missing block is a staircase. Neither alone.
  \end{itemize}
\item
  Matrix completion accommodates all three and so \textbf{unites}
  synthetic control and unconfoundedness in one estimator. Both of their
  assumptions reappear in its identification.
\end{itemize}
\end{frame}

\begin{frame}{Identification, estimation, inference}
\protect\phantomsection\label{identification-estimation-inference}
\begin{itemize}
\tightlist
\item
  Model: \(\mathbf Y = \mathbf L^{*} + \boldsymbol\varepsilon\),
  \(\mathbb E[\boldsymbol\varepsilon \mid \mathbf L^{*}] = 0\), errors
  sub-Gaussian and independent. Two assumptions:

  \begin{enumerate}
  \tightlist
  \item
    \textbf{Unconfoundedness given \(\mathbf L^{*}\).} Adoption dates
    are independent of the errors, conditional on the systematic
    component. Timing may depend on the unit's latent type, not on luck.
  \item
    \textbf{Low rank.} \(\mathbf L^{*}\) is driven by a few latent
    factors; rows and columns are nearly linearly dependent. With many
    independent dimensions there would be too little to borrow.
  \end{enumerate}
\item
  Estimator: minimise squared error on the \textbf{observed} cells, with
  unit and time fixed effects and a penalty on the \textbf{nuclear norm}
  (the sum of singular values, the convex relaxation of rank): \[
  \min_{\mathbf L, \Gamma, \Delta}\ \frac{1}{|\mathcal O|} \big\| P_{\mathcal O}\big( \mathbf Y - \mathbf L - \Gamma \mathbf 1_T' - \mathbf 1_N \Delta' \big) \big\|_F^2 + \Lambda \| \mathbf L \|_{*}.
  \] \(P_{\mathcal O}\) keeps the observed entries; \(\Lambda\) by
  cross-validation on held-out observed cells.
\item
  Algorithm (soft-impute): fill the missing cells with the current
  \(\hat{\mathbf L}\), remove the fixed effects, take the SVD,
  \textbf{shrink every singular value by \(\Lambda\)} and drop those
  that hit zero, repeat. Inference: a nonparametric bootstrap over
  units.
\end{itemize}
\end{frame}

\begin{frame}{Proposition 99 by matrix completion}
\protect\phantomsection\label{proposition-99-by-matrix-completion}
\begin{center}
\includegraphics[width=0.72\linewidth,height=\textheight,keepaspectratio]{08_synthetic_control_files/figure-beamer/mc-fig-1.pdf}
\end{center}

\begin{center}\scriptsize
\begin{tabular}{lrrrrrrr}
\toprule
$\Lambda$ & 2 & 5 & 10 & 20 & 40 & 80 & 160 \\
\midrule
Rank of $\hat{\mathbf L}$ &  30 & 25 & 19 & 11 & 5 & 4 & 1  \\
Pre-period RMSPE &  0.23 & 0.54 & 0.96 & 1.62 & 2.55 & 4.11 & 5.74  \\
Mean gap 1989--2000 &  -19.7 & -19.8 & -19.9 & -20.0 & -20.2 & -21.0 & -23.7  \\
\bottomrule
\end{tabular}
\end{center}

\begin{itemize}
\tightlist
\item
  The \(39 \times 31\) sales matrix with California's 12 post-period
  cells blanked; two-way fixed effects plus a shrunk low-rank component,
  by hand. Cross-validation picks \(\Lambda = 10\), rank 19, mean gap
  \(-19.9\).
\item
  The estimate barely moves across ranks from 30 down to 4. The method
  does not need the non-negativity or the convex hull; it needs the
  factors to be few and the timing to be unrelated to the noise.
\end{itemize}
\end{frame}

\section{7. Synthetic
Difference-in-Differences}\label{synthetic-difference-in-differences}

\begin{frame}{Dylan and Cash}
\protect\phantomsection\label{dylan-and-cash}
\begin{itemize}
\tightlist
\item
  Synthetic difference-in-differences (Arkhangelsky, Athey, Hirshberg,
  Imbens and Wager 2021) is a duet. From \textbf{synthetic control}:
  unit weights that make the donors track the treated unit before
  treatment. From \textbf{DiD}: unit and time fixed effects, so a level
  gap between the two series is allowed and absorbed. New to both:
  \textbf{time weights}, which pick the pre-periods most informative
  about the post-period.
\item
  Why: synthetic control is biased when the fit is imperfect, DiD when
  trends are not parallel. SDID needs neither the convex hull nor
  parallel trends in the raw data. It \textbf{approximates parallel
  trends by reweighting}, then runs a DiD on the reweighted panel.
  Demeaned synthetic control allowed an intercept; SDID adds regularised
  unit weights and time weights.
\end{itemize}

\vspace{0.4em}

\begin{itemize}
\tightlist
\item
  Outcome model, close to matrix completion's but without the low-rank
  penalty: \[
  \mathbf Y = \mathbf L + \mathbf D \circ \boldsymbol\delta + \mathbf E,
  \] \(\mathbf L = \Gamma \Upsilon'\) the systematic component (unit
  factors times time factors), \(\mathbf D \circ \boldsymbol\delta\) the
  effect applied only where \(D_{it} = 1\), \(\mathbf E\) noise.
  \(\mathbf L\) is where the bias lives, so \(\mathbf L\) is what the
  weights balance.
\end{itemize}
\end{frame}

\begin{frame}{Three steps}
\protect\phantomsection\label{three-steps}
\begin{enumerate}
\tightlist
\item
  \textbf{Unit weights} \(\hat\omega\): a synthetic control fit with an
  intercept and a ridge penalty,
  \(\hat\omega_0 + \hat\omega' \mathbf Y_{\text{co},\text{pre}} \approx \bar{\mathbf Y}_{\text{tr},\text{pre}}\):
  \[
  (\hat\omega_0, \hat\omega) = \arg\min_{\omega_0 \in \mathbb R,\ \omega \in \Omega}\ \textstyle\sum_{t \leq T_0} \big( \omega_0 + \sum_{i \in \text{co}} \omega_i Y_{it} - \bar Y_{\text{tr},t} \big)^2 + \zeta^2 T_0 \|\omega\|_2^2 .
  \] \(\Omega\) is the simplex;
  \(\zeta = (N_{\text{tr}} T_{\text{post}})^{1/4} \hat\sigma\), with
  \(\hat\sigma\) the typical one-period change among controls, keeps the
  weights from chasing noise.
\item
  \textbf{Time weights} \(\hat\lambda\): the same problem transposed.
  Weight pre-periods so that, for the controls, the weighted pre-period
  average predicts the post-period average, up to an intercept: \[
  (\hat\lambda_0, \hat\lambda) = \arg\min_{\lambda_0,\ \lambda \in \Lambda}\ \textstyle\sum_{i \in \text{co}} \big( \lambda_0 + \sum_{t \leq T_0} \lambda_t Y_{it} - \bar Y_{i,\text{post}} \big)^2 .
  \]
\item
  \textbf{A weighted two-way fixed-effects regression}, with
  \(\hat\omega_i = 1/N_{\text{tr}}\) for treated units and
  \(\hat\lambda_t = 1/T_{\text{post}}\) for post-periods: \[
  \hat\tau^{\text{sdid}} = \arg\min_{\tau, \mu, \alpha, \beta} \textstyle\sum_{i=1}^{N} \sum_{t=1}^{T} \big( Y_{it} - \mu - \alpha_i - \beta_t - D_{it}\tau \big)^2\, \hat\omega_i \hat\lambda_t .
  \] With one treated unit this is four weighted averages and three
  subtractions.
\end{enumerate}
\end{frame}

\begin{frame}{Four assumptions}
\protect\phantomsection\label{four-assumptions}
\begin{enumerate}
\tightlist
\item
  \textbf{Errors.} Rows of \(\mathbf E\) are independent Gaussian
  vectors with bounded covariance eigenvalues. Stronger than DiD or
  synthetic control, which assume independence but not a distribution.
\item
  \textbf{Sample sizes.} Many controls \(N_{\text{co}}\) and many
  pre-periods \(T_{\text{pre}}\), in balance; ideally
  \(N_{\text{tr}} T_{\text{post}}\) large as well. The method needs room
  to reweight in both directions.
\item
  \textbf{Systematic component.} \(\mathbf L\) has a few large singular
  values: low rank again, though not imposed by a penalty.
\item
  \textbf{Oracle weights.} There exist weights, computable from
  \(\mathbf L\), that remove the systematic difference between treated
  and control; the sample weights estimate them. This is
  \textbf{parallel trends after reweighting}: not parallel in the raw
  data, parallel once units and periods are reweighted, and only
  parallel, not identical.
\end{enumerate}

\vspace{0.4em}

\begin{itemize}
\tightlist
\item
  The estimated weights converge to the oracle weights as the panel
  grows in both dimensions; that is what makes the estimator consistent.
\item
  DiD (uniform weights, parallel trends in the raw data), synthetic
  control (unit weights, no intercept, no time weights, exact fit) and
  SDID are three points on one line.
\end{itemize}
\end{frame}

\begin{frame}{Proposition 99 by synthetic DiD}
\protect\phantomsection\label{proposition-99-by-synthetic-did}
\begin{center}
\includegraphics[width=0.72\linewidth,height=\textheight,keepaspectratio]{08_synthetic_control_files/figure-beamer/sdid-fig-1.pdf}
\end{center}

\begin{center}\scriptsize
\begin{tabular}{lrrr}
\toprule
 & DiD & Synthetic control & Synthetic DiD \\
\midrule
Estimate & -27.3 & -19.5 & -15.6 \\
Placebo standard error & 19.4 & 10.2 & 10.2 \\
Arkhangelsky et al.\ (2021), Table 1 & $-27.3$ & $-19.6$ & $-15.6$ \\
\bottomrule
\end{tabular}
\end{center}

\begin{itemize}
\tightlist
\item
  Unit weights spread over 22 states (largest Nevada 0.12); the control
  sits 24.7 packs above California before 1989 and the intercept absorbs
  it. Time weights: 1986 0.37, 1987 0.21, 1988 0.43; the 1970s get
  nothing.
\end{itemize}
\end{frame}

\begin{frame}[fragile]{Practical considerations}
\protect\phantomsection\label{practical-considerations}
\begin{itemize}
\tightlist
\item
  \textbf{Look at the pre-trends.} SDID lets the treated unit and its
  control sit apart by an intercept, so the two series will not lie on
  top of each other as a good synthetic control does. Check that they
  are parallel after weighting; visual inspection is still the main
  tool.
\item
  \textbf{Show the weights.} Time weights under the time-series plot, as
  on the previous slide, tell the reader which pre-years drive the
  counterfactual; unit weights tell them who is in it. The weights are
  what let parallel trends fail in the raw data and hold in the
  estimate, so they are the design.
\item
  \textbf{Show all three.} DiD, synthetic control and SDID side by side,
  as Arkhangelsky et al.~do. Disagreement tells you which assumption is
  doing the work.
\item
  \textbf{Staggered adoption} is covered: the original appendix, section
  2.3 of Clarke, Pailañir, Athey and Imbens (2024), and Porreca (2022)
  extend the estimator cohort by cohort.
\item
  \textbf{Inference.} Placebo (as on the previous slide), jackknife or
  bootstrap. With few treated units, placebo; with many, bootstrap or
  jackknife over units.
\item
  Software: \texttt{synthdid} in R, \texttt{sdid} in Stata. This deck
  solved both weight problems with a quadratic programme.
\end{itemize}
\end{frame}

\section{8. Using It}\label{using-it}

\begin{frame}{When synthetic control is the design}
\protect\phantomsection\label{when-synthetic-control-is-the-design}
\begin{itemize}
\tightlist
\item
  The setting: \textbf{one or a few aggregate treated units}, a long
  pre-period, a pool of untreated aggregates. A state law, a country's
  capital-account reform, one exchange's tick-size change, one
  currency's peg. DiD with hand-picked controls is Card's design;
  synthetic control chooses the controls by fit and does inference by
  placebo.
\item
  Abadie (2021) lists what the data must offer before the method is
  credible:

  \begin{enumerate}
  \tightlist
  \item
    an effect \textbf{large relative to the volatility} of the outcome,
    since one unit's noise cannot be averaged away;
  \item
    a donor pool \textbf{not affected} by the treatment, by similar
    interventions, or by spillovers;
  \item
    \textbf{no anticipation}, or a treatment date moved back to when
    behaviour changed;
  \item
    the treated unit \textbf{inside the convex hull}: not the largest,
    smallest or strangest unit in the pool;
  \item
    a \textbf{long pre-period} and enough post-periods for the effect to
    appear.
  \end{enumerate}
\item
  Examples that follow the recipe: Billmeier and Nannicini (2013) on
  economic liberalisation; Born et al.~(2019) on the Brexit vote and UK
  GDP; Cavallo et al.~(2013) on disasters and growth. With hundreds of
  treated firms you are back in Lecture 7b, or in SDID and matrix
  completion with staggered timing.
\end{itemize}
\end{frame}

\begin{frame}{Conclusion}
\protect\phantomsection\label{conclusion}
\begin{enumerate}
\tightlist
\item
  \textbf{Synthetic control} imputes one treated unit's untreated
  outcome from a non-negative, sum-to-one weighted average of donors,
  chosen by pre-period fit. It assumes a factor model, not parallel
  trends; the long, close fit is the evidence.
\item
  \textbf{Inference is by placebo.} Each donor as if treated;
  post-to-pre RMSPE ratios; report the rank. A placebo in time checks
  the date.
\item
  \textbf{Non-negative weights are interpolation.} They refuse to fit
  outliers, which keeps the estimate free of model dependence.
  Regression fits better by extrapolating.
\item
  \textbf{Imperfect fit is bias.} Demeaning allows a level gap;
  augmentation corrects the matching bias with a ridge model, at the
  price of negative weights. Report fit, weights and several
  specifications.
\item
  \textbf{Matrix completion} fills the potential-outcome matrix under
  low rank and unconfoundedness given the factors.
\item
  \textbf{Synthetic DiD} reweights units and periods, allows an
  intercept, runs a weighted DiD: parallel trends after reweighting.
  Show DiD, SC and SDID together, and the weights.
\item
  \textbf{Which units are inside my counterfactual, with what weight,
  and what did I give up to fit better?}
\end{enumerate}
\end{frame}

\begin{frame}{References}
\protect\phantomsection\label{references}
\tiny

\begin{itemize}
\tightlist
\item
  Abadie, A. (2021). Using Synthetic Controls: Feasibility, Data
  Requirements, and Methodological Aspects. \emph{Journal of Economic
  Literature} 59(2), 391--425.
\item
  Abadie, A., A. Diamond and J. Hainmueller (2010). Synthetic Control
  Methods for Comparative Case Studies: Estimating the Effect of
  California's Tobacco Control Program. \emph{JASA} 105(490), 493--505.
\item
  Abadie, A., A. Diamond and J. Hainmueller (2015). Comparative Politics
  and the Synthetic Control Method. \emph{American Journal of Political
  Science} 59(2), 495--510.
\item
  Abadie, A. and J. Gardeazabal (2003). The Economic Costs of Conflict:
  A Case Study of the Basque Country. \emph{AER} 93(1), 113--132.
\item
  Abadie, A. and G. Imbens (2011). Bias-Corrected Matching Estimators
  for Average Treatment Effects. \emph{JBES} 29(1), 1--11.
\item
  Arkhangelsky, D., S. Athey, D. Hirshberg, G. Imbens and S. Wager
  (2021). Synthetic Difference-in-Differences. \emph{AER} 111(12),
  4088--4118.
\item
  Athey, S., M. Bayati, N. Doudchenko, G. Imbens and K. Khosravi (2021).
  Matrix Completion Methods for Causal Panel Data Models. \emph{JASA}
  116(536), 1716--1730.
\item
  Ben-Michael, E., A. Feller and J. Rothstein (2021). The Augmented
  Synthetic Control Method. \emph{JASA} 116(536), 1789--1803.
\item
  Billmeier, A. and T. Nannicini (2013). Assessing Economic
  Liberalization Episodes: A Synthetic Control Approach. \emph{REStat}
  95(3), 983--1001.
\item
  Born, B., G. Müller, M. Schularick and P. Sebastian (2019). The Costs
  of Economic Nationalism: Evidence from the Brexit Experiment.
  \emph{Economic Journal} 129(623), 2722--2744.
\item
  Card, D. (1990). The Impact of the Mariel Boatlift on the Miami Labor
  Market. \emph{ILR Review} 43(2), 245--257.
\item
  Cavallo, E., S. Galiani, I. Noy and J. Pantano (2013). Catastrophic
  Natural Disasters and Economic Growth. \emph{REStat} 95(5),
  1549--1561.
\item
  Clarke, D., D. Pailañir, S. Athey and G. Imbens (2024). On Synthetic
  Difference-in-Differences and Related Estimation Methods in Stata.
  \emph{Stata Journal} 24(4), 557--598.
\item
  Cunningham, S. (2025). \emph{Causal Inference: The Remix}, ch.~11 ---
  the source for this lecture and for both datasets.
\item
  Doudchenko, N. and G. Imbens (2017). Balancing, Regression,
  Difference-in-Differences and Synthetic Control Methods: A Synthesis.
  Working paper.
\item
  Ferman, B. and C. Pinto (2021). Synthetic Controls with Imperfect
  Pretreatment Fit. \emph{Quantitative Economics} 12(4), 1197--1221.
\item
  Ferman, B., C. Pinto and V. Possebom (2020). Cherry Picking with
  Synthetic Controls. \emph{Journal of Policy Analysis and Management}
  39(2), 510--532.
\item
  Kaul, A., S. Klößner, G. Pfeifer and M. Schieler (2022). Standard
  Synthetic Control Methods: The Case of Using All Preintervention
  Outcomes Together with Covariates. \emph{JBES} 40(3), 1362--1376.
\item
  King, G. and L. Zeng (2006). The Dangers of Extreme Counterfactuals.
  \emph{Political Analysis} 14(2), 131--159.
\item
  Peri, G. and V. Yasenov (2019). The Labor Market Effects of a Refugee
  Wave: Synthetic Control Method Meets the Mariel Boatlift.
  \emph{Journal of Human Resources} 54(2), 267--309.
\end{itemize}
\end{frame}

\section{Backup}\label{backup}

\begin{frame}[fragile]{Stata and R: the estimators of this lecture}
\protect\phantomsection\label{stata-and-r-the-estimators-of-this-lecture}
\begin{Shaded}
\begin{Highlighting}[]
\KeywordTok{use}\NormalTok{ https:}\CommentTok{//github.com/scunning1975/mixtape/raw/master/smoking.dta, clear}
\KeywordTok{tsset}\NormalTok{ state }\FunctionTok{year}
\NormalTok{synth cigsale lnincome(1980(1)1988) age15to24(1980(1)1988) retprice(1980(1)1988) }\CommentTok{///}
\NormalTok{    beer(1984(1)1988) cigsale(1988) cigsale(1980) cigsale(1975), }\CommentTok{///}
\NormalTok{    trunit(3) trperiod(1989) unitnames(state) fig            }\CommentTok{// ADH (2010)}
\NormalTok{synth\_runner cigsale ..., trunit(3) trperiod(1989) gen\_vars  }\CommentTok{// placebos, ratios, p{-}values}
\NormalTok{allsynth cigsale cigsale(1970(1)1988), trunit(3) trperiod(1989) }\CommentTok{///}
\NormalTok{    bcorrect(}\KeywordTok{merge}\NormalTok{ ridge figure) gapfigure(classic bcorrect)    }\CommentTok{// augmented SC (Wiltshire)}
\NormalTok{sdid cigsale state }\FunctionTok{year}\NormalTok{ treated, }\KeywordTok{vce}\NormalTok{(placebo) reps(200) }\KeywordTok{graph}  \CommentTok{// synthetic DiD (Clarke et al.)}
\end{Highlighting}
\end{Shaded}

\begin{itemize}
\tightlist
\item
  R: \texttt{Synth} (\texttt{dataprep}, \texttt{synth}),
  \texttt{augsynth}, \texttt{MCPanel} (\texttt{mcnnm\_cv}),
  \texttt{synthdid}, \texttt{gsynth}.
\item
  This deck: the inner problem is \texttt{quadprog::solve.QP}; the \(V\)
  search is \texttt{optim}; augmentation is one ridge solve; matrix
  completion is an SVD soft-threshold loop; SDID is two quadratic
  programmes and four weighted averages. No synthetic-control package
  was loaded.
\end{itemize}
\end{frame}




\end{document}
